\begin{abstract}

    This work investigates two least-squares problems using data derived
    from the student's enrollment number. First, an overdetermined linear
    system $A\mathbf{x}\approx\mathbf{b}$, with
    $A\in\mathbb{R}^{5\times4}$, is solved using the QR factorization
    $A=QR$. The least-squares solution is obtained by solving the
    resulting upper triangular system. The existence of a unique solution
    is discussed in terms of the rank of the coefficient matrix.

    Second, a quadratic polynomial is fitted to six experimental data
    points using the least-squares formulation. The associated design
    matrix is constructed, and the polynomial coefficients are determined
    computationally. The quality of the fit is assessed by plotting the
    experimental data and the fitted curve and calculating the mean
    squared error.

    The numerical implementation explicitly constructs the principal
    linear algebra algorithms, avoiding high-level routines for QR
    factorization and least-squares solutions. This approach connects
    the mathematical foundations of least squares with their
    computational implementation and numerical evaluation.

\end{abstract}

\noindent\textbf{Keywords:} QR factorization; least squares;
quadratic polynomial fitting; numerical linear algebra.